% =============================================================================
% estilo.tex — Identidad visual del informe ("papel y tinta", igual que la app)
% =============================================================================
\usepackage[letterpaper,top=2.5cm,bottom=2.5cm,left=2.6cm,right=2.4cm,headheight=14pt]{geometry}
\usepackage{mathpazo}                 % matemáticas Palatino (antes de fontspec)
\usepackage[no-math]{fontspec}
\usepackage{babel}
\babelprovide[import,main]{spanish}
\setlocalecaption{spanish}{table}{Tabla}
\setlocalecaption{spanish}{contents}{Contenido}
\setmainfont{TeX Gyre Pagella}
\newfontfamily\titular{Fraunces Proj}[UprightFont=* SemiBold,BoldFont=* Black,ItalicFont=* Italic,
  BoldItalicFont=* Bold Italic]
\newfontfamily\titularfino{Fraunces Proj}[UprightFont=* Regular,ItalicFont=* Italic]
\newfontfamily\sans{Manrope Proj}[UprightFont=* Regular,BoldFont=* Bold]
\newfontfamily\sansmedio{Manrope Proj}[UprightFont=* SemiBold,BoldFont=* ExtraBold]
\setmonofont{JetBrains Mono}[Scale=0.86]
\newfontfamily\mono{JetBrains Mono}[Scale=0.86]

\usepackage{microtype}
\usepackage{xcolor}
\definecolor{papel}{HTML}{F4EFE6}
\definecolor{tarjeta}{HTML}{FFFDF8}
\definecolor{tinta}{HTML}{1D1A16}
\definecolor{gris}{HTML}{6E665B}
\definecolor{linea}{HTML}{DCD2C1}
\definecolor{bermellon}{HTML}{D9480F}
\definecolor{oliva}{HTML}{5E7D3A}
\definecolor{pizarra}{HTML}{3E5C76}
\definecolor{crema}{HTML}{FBF4E9}
\definecolor{gristeoria}{HTML}{F0EFEA}

\usepackage{amsmath,amssymb,mathtools}
\usepackage{graphicx}
\usepackage{booktabs,tabularx,array,multirow}
\usepackage{caption,subcaption}
\usepackage{enumitem}
\usepackage{fancyhdr}
\usepackage{titlesec}
\usepackage{tikz}
\usetikzlibrary{calc,positioning,shadows}
\usepackage[most]{tcolorbox}
\usepackage{eso-pic}
\usepackage{float}
\usepackage{setspace}
\usepackage{needspace}
\usepackage[hidelinks]{hyperref}
\hypersetup{colorlinks=true,linkcolor=tinta,urlcolor=bermellon,citecolor=bermellon}
\usepackage{bookmark}

\setstretch{1.08}
\setlength{\parskip}{0.42em}
\setlength{\parindent}{0pt}
\renewcommand{\arraystretch}{1.12}
\captionsetup{font={small,sf},labelfont={bf,color=bermellon},textfont={color=tinta},labelsep=period,
  skip=4pt}
\DeclareCaptionFont{sf}{\sans}
\setlist{itemsep=2pt,topsep=3pt,leftmargin=1.3em}

% ---------- Títulos de sección ------------------------------------------------
\titleformat{\section}[block]
  {\titular\fontsize{19}{22}\selectfont\color{tinta}}
  {\color{bermellon}\titular\bfseries\fontsize{30}{30}\selectfont\thesection\hspace{10pt}}{0pt}{}
  [\vspace{-4pt}{\color{tinta}\rule{\linewidth}{1.4pt}}]
\titleformat{\subsection}{\titular\fontsize{13.2}{16}\selectfont\color{tinta}}
  {\color{bermellon}\thesubsection}{8pt}{}
\titleformat{\subsubsection}{\sansmedio\fontsize{10.4}{13}\selectfont\color{pizarra}}
  {\thesubsubsection}{6pt}{}
\titlespacing*{\section}{0pt}{10pt}{8pt}
\pretocmd{\section}{\needspace{14\baselineskip}}{}{}
\pretocmd{\subsection}{\needspace{8\baselineskip}}{}{}
\titlespacing*{\subsection}{0pt}{9pt}{3pt}
\titlespacing*{\subsubsection}{0pt}{6pt}{2pt}

% ---------- Encabezados y pies -------------------------------------------------
\pagestyle{fancy}
\fancyhf{}
\renewcommand{\headrulewidth}{0pt}
\fancyhead[L]{\mono\color{gris}\fontsize{7}{8}\selectfont ESTADÍSTICA PARA INGENIERÍA · PROYECTO DE AULA 2026-2}
\fancyhead[R]{\mono\color{gris}\fontsize{7}{8}\selectfont UNIVERSIDAD DE MEDELLÍN}
\fancyfoot[C]{\sans\color{tinta}\small\thepage}
\fancypagestyle{plain}{\fancyhf{}\fancyfoot[C]{\sans\color{tinta}\small\thepage}}

% ---------- Cajas ---------------------------------------------------------------
\newcommand{\etiquetacaja}[2]{{\mono\color{#1}\fontsize{7}{8}\selectfont #2}\par\vspace{2pt}}
\newtcolorbox{lectura}[1][]{enhanced,breakable,colback=crema,colframe=bermellon,boxrule=0pt,
  leftrule=3pt,arc=0pt,left=9pt,right=8pt,top=5pt,bottom=5pt,fontupper=\small,
  before upper=\etiquetacaja{bermellon}{LECTURA},#1}
\newtcolorbox{idea}[1][]{enhanced,breakable,colback=gristeoria,colframe=pizarra,boxrule=0pt,
  leftrule=3pt,arc=0pt,left=9pt,right=8pt,top=5pt,bottom=5pt,fontupper=\small,
  before upper=\etiquetacaja{pizarra}{#1}}
\newtcbtheorem[number within=section]{definicion}{Definición}{enhanced,breakable,colback=white,
  colframe=tinta,boxrule=0.9pt,arc=0pt,fonttitle=\sansmedio\small,coltitle=papel,colbacktitle=tinta,
  attach boxed title to top left={yshift=-2mm,xshift=3mm},boxed title style={arc=0pt,boxrule=0pt},
  top=4mm,left=6pt,right=6pt,bottom=3pt,fontupper=\small,
  drop fuzzy shadow=black!0,borderline west={3pt}{0pt}{bermellon}}{def}
\newtcbtheorem[number within=section]{teorema}{Teorema}{enhanced,breakable,colback=crema,
  colframe=bermellon,boxrule=0.9pt,arc=0pt,fonttitle=\sansmedio\small,coltitle=white,colbacktitle=bermellon,
  attach boxed title to top left={yshift=-2mm,xshift=3mm},boxed title style={arc=0pt,boxrule=0pt},
  top=4mm,left=6pt,right=6pt,bottom=3pt,fontupper=\small\itshape}{teo}
\newtcolorbox{algoritmo}[1]{enhanced,breakable,colback=white,colframe=pizarra,boxrule=0.8pt,arc=0pt,
  fonttitle=\sansmedio\small,coltitle=white,colbacktitle=pizarra,title={#1},left=6pt,right=6pt,
  fontupper=\small}

% ---------- Utilidades -----------------------------------------------------------
\newcommand{\etiqueta}[1]{{\mono\color{bermellon}\fontsize{7.5}{9}\selectfont\MakeUppercase{#1}}}
\newcommand{\rechaza}{\colorbox{bermellon}{\color{white}\sansmedio\footnotesize\,SE RECHAZA $H_0$\,}}
\newcommand{\norechaza}{\colorbox{oliva}{\color{white}\sansmedio\footnotesize\,NO SE RECHAZA $H_0$\,}}
\newcommand{\figura}[3][0.98]{\begin{figure}[H]\centering\includegraphics[width=#1\linewidth]{../figuras/#2.pdf}
  \caption{#3}\label{fig:#2}\end{figure}}
\newcommand{\tabla}[3][\small]{\begin{table}[H]\centering\caption{#3}\label{tab:#2}{#1\sans
  \input{tablas/#2.tex}}\end{table}}
\newcommand{\pyval}[1]{#1}
